% GENERATED FILE. Edit canonical lessons, then run npm run generate:books.
% canonical-source-hash: fnv1a64:f79efac997c6b969
% canonical-lessons: JA-C208-hare, JA-C208-kumori, JA-C208-tenkiyohou, JA-C208-sora, JA-C208-taifuu

\chapter{The Weather Forecast}
\label{ch:ja-the-weather-forecast}
\hlchaptermodality{208}

\begin{chapteropening}
\textbf{By the end of this chapter:} \emph{I can understand a simple weather forecast.}

Complete the last lesson of chapter 208: I can understand a simple weather forecast.
\end{chapteropening}

\section[\texorpdfstring{\ja{はれ} (hare)}{はれ (hare)}]{\ja{はれ} (hare) — fine weather; clear skies}

\label{lesson:JA-C208-hare}



\begin{quote}
Before the new one: say the Japanese for to clear up, then the Japanese for to cloud over.
\end{quote}

\subsection*{You'll want to know: \ja{はれ}}
\textbf{\ja{はれ}} — \emph{hare} — \textquotedblleft{}fine weather; clear skies\textquotedblright{}.

\begin{quote}
\textbf{\ja{あしたは} \ja{はれでしょう。}} — \emph{ashita wa hare deshō} — \textquotedblleft{}Tomorrow will probably be fine.\textquotedblright{}
\end{quote}

\begin{cousinweb}
\textbf{\ja{はれ}} (\emph{hare}) is native Japanese (wago), derived from the verb \emph{hareru}, \textquotedblleft{}to clear up\textquotedblright{}, inherited rather than borrowed (JMdict).
\end{cousinweb}

\begin{grammarlens}[title={What you've built}]
One more word for everyday life at A2.
\end{grammarlens}

\subsection*{Your turn}
\begin{itemize}
\raggedright
  \item \emph{Say it:} \emph{hare}
  \item \emph{Say it:} \emph{hare}, once more
  \item \emph{Say it:} \emph{ashita wa hare deshō} — Tomorrow will probably be fine.
  \item \emph{Recall:} say the Japanese for to clear up, then the Japanese for to cloud over, then say \emph{hare} again
\end{itemize}

\begin{tcolorbox}[breakable,colback=teal!4,colframe=teal!35!black,title={Before you move on}]
What does \ja{はれ} mean? (\textquotedblleft{}Fine weather; clear skies\textquotedblright{}.) Say it once more.
\end{tcolorbox}
\section[\texorpdfstring{\ja{くもり} (kumori)}{くもり (kumori)}]{\ja{くもり} (kumori) — cloudy weather}

\label{lesson:JA-C208-kumori}



\begin{quote}
Before the new one: say the Japanese for to cloud over, then the Japanese for fine weather.
\end{quote}

\subsection*{You'll want to know: \ja{くもり}}
\textbf{\ja{くもり}} — \emph{kumori} — \textquotedblleft{}cloudy weather\textquotedblright{}.

\begin{quote}
\textbf{\ja{らいしゅうは} \ja{くもりでしょう。}} — \emph{raishū wa kumori deshō} — \textquotedblleft{}Next week will probably be cloudy.\textquotedblright{}
\end{quote}

\begin{cousinweb}
\textbf{\ja{くもり}} (\emph{kumori}) is native Japanese (wago), derived from the verb \emph{kumoru}, \textquotedblleft{}to cloud over\textquotedblright{}, inherited rather than borrowed (JMdict).
\end{cousinweb}

\begin{grammarlens}[title={What you've built}]
One more word for everyday life at A2.
\end{grammarlens}

\subsection*{Your turn}
\begin{itemize}
\raggedright
  \item \emph{Say it:} \emph{kumori}
  \item \emph{Say it:} \emph{kumori}, once more
  \item \emph{Say it:} \emph{raishū wa kumori deshō} — Next week will probably be cloudy.
  \item \emph{Recall:} say the Japanese for to cloud over, then the Japanese for fine weather, then say \emph{kumori} again
\end{itemize}

\begin{tcolorbox}[breakable,colback=teal!4,colframe=teal!35!black,title={Before you move on}]
What does \ja{くもり} mean? (\textquotedblleft{}Cloudy weather\textquotedblright{}.) Say it once more.
\end{tcolorbox}
\section[\texorpdfstring{\ja{てんきよほう} (tenkiyohō)}{てんきよほう (tenkiyohō)}]{\ja{てんきよほう} (tenkiyohō) — a weather forecast}

\label{lesson:JA-C208-tenkiyohou}



\begin{quote}
Before the new one: say the Japanese for fine weather, then the Japanese for cloudy weather.
\end{quote}

\subsection*{You'll want to know: \ja{てんきよほう}}
\textbf{\ja{てんきよほう}} — \emph{tenkiyohō} — \textquotedblleft{}a weather forecast\textquotedblright{}.

\begin{quote}
\textbf{\ja{あしたの} \ja{てんきよほうは} \ja{ゆきです。}} — \emph{ashita no tenkiyohō wa yuki desu} — \textquotedblleft{}Tomorrow's forecast is snow.\textquotedblright{}
\end{quote}

\begin{cousinweb}
\textbf{\ja{てんきよほう}} (\emph{tenkiyohō}) is Sino-Japanese (kango): four Chinese characters read with their on'yomi, readings borrowed from Chinese. They are \emph{ten}, \textquotedblleft{}sky\textquotedblright{} (Middle Chinese \emph{then}), \emph{ki}, \textquotedblleft{}air\textquotedblright{} (\emph{khj+jH}), \emph{yo}, \textquotedblleft{}beforehand\textquotedblright{} (\emph{yoH}) and \emph{hō}, \textquotedblleft{}report\textquotedblright{} (\emph{pawH}) (KANJIDIC2; Middle Chinese in Baxter's transcription).
\end{cousinweb}

\begin{grammarlens}[title={What you've built}]
One more word for everyday life at A2.
\end{grammarlens}

\subsection*{Your turn}
\begin{itemize}
\raggedright
  \item \emph{Say it:} \emph{tenkiyohō}
  \item \emph{Say it:} \emph{tenkiyohō}, once more
  \item \emph{Say it:} \emph{ashita no tenkiyohō wa yuki desu} — Tomorrow's forecast is snow.
  \item \emph{Recall:} say the Japanese for fine weather, then the Japanese for cloudy weather, then say \emph{tenkiyohō} again
\end{itemize}

\begin{tcolorbox}[breakable,colback=teal!4,colframe=teal!35!black,title={Before you move on}]
What does \ja{てんきよほう} mean? (\textquotedblleft{}A weather forecast\textquotedblright{}.) Say it once more.
\end{tcolorbox}
\section[\texorpdfstring{\ja{そら} (sora)}{そら (sora)}]{\ja{そら} (sora) — the sky}

\label{lesson:JA-C208-sora}



\begin{quote}
Before the new one: say the Japanese for cloudy weather, then the Japanese for a weather forecast.
\end{quote}

\subsection*{You'll want to know: \ja{そら}}
\textbf{\ja{そら}} — \emph{sora} — \textquotedblleft{}the sky\textquotedblright{}.

\begin{quote}
\textbf{\ja{あしたは} \ja{そらが} \ja{あおいでしょう。}} — \emph{ashita wa sora ga aoi deshō} — \textquotedblleft{}The sky will probably be blue tomorrow.\textquotedblright{}
\end{quote}

\begin{cousinweb}
\textbf{\ja{そら}} (\emph{sora}) is a native Japanese word (wago), inherited rather than borrowed from Chinese; its kanji was attached later as a kun'yomi, and its older origin is not settled (JMdict, KANJIDIC2).
\end{cousinweb}

\begin{grammarlens}[title={What you've built}]
One more word for everyday life at A2.
\end{grammarlens}

\subsection*{Your turn}
\begin{itemize}
\raggedright
  \item \emph{Say it:} \emph{sora}
  \item \emph{Say it:} \emph{sora}, once more
  \item \emph{Say it:} \emph{ashita wa sora ga aoi deshō} — The sky will probably be blue tomorrow.
  \item \emph{Recall:} say the Japanese for cloudy weather, then the Japanese for a weather forecast, then say \emph{sora} again
\end{itemize}

\begin{tcolorbox}[breakable,colback=teal!4,colframe=teal!35!black,title={Before you move on}]
What does \ja{そら} mean? (\textquotedblleft{}The sky\textquotedblright{}.) Say it once more.
\end{tcolorbox}
\section[\texorpdfstring{\ja{たいふう} (taifū)}{たいふう (taifū)}]{\ja{たいふう} (taifū) — a typhoon}

\label{lesson:JA-C208-taifuu}



\begin{quote}
Before the new one: say the Japanese for a weather forecast, then the Japanese for the sky.
\end{quote}

\subsection*{You'll want to know: \ja{たいふう}}
\textbf{\ja{たいふう}} — \emph{taifū} — \textquotedblleft{}a typhoon\textquotedblright{}.

\begin{quote}
\textbf{\ja{たいふうは} \ja{らいしゅうでしょう。}} — \emph{taifū wa raishū deshō} — \textquotedblleft{}The typhoon will probably come next week.\textquotedblright{}
\end{quote}

\begin{cousinweb}
\textbf{\ja{たいふう}} (\emph{taifū}) is Sino-Japanese (kango), read with on'yomi borrowed from Chinese. Its second part is \emph{fū}, \textquotedblleft{}wind\textquotedblright{} (Middle Chinese \emph{pjuwng}); how the first part, \emph{tai}, came to name a great storm is not settled (JMdict, KANJIDIC2; Baxter's transcription).
\end{cousinweb}

\begin{grammarlens}[title={What you've built}]
One more word for everyday life at A2.
\end{grammarlens}

\subsection*{Your turn}
\begin{itemize}
\raggedright
  \item \emph{Say it:} \emph{taifū}
  \item \emph{Say it:} \emph{taifū}, once more
  \item \emph{Say it:} \emph{taifū wa raishū deshō} — The typhoon will probably come next week.
  \item \emph{Recall:} say the Japanese for a weather forecast, then the Japanese for the sky, then say \emph{taifū} again
\end{itemize}

\begin{tcolorbox}[breakable,colback=teal!4,colframe=teal!35!black,title={Before you move on}]
What does \ja{たいふう} mean? (\textquotedblleft{}A typhoon\textquotedblright{}.) Say it once more.
\end{tcolorbox}
